\setchapterpreamble[u]{\margintoc}
\chapter{Регулярные языки}
\label{chpt:RegularLanguages}
\tikzsetfigurename{RegularLanguages_}

Регулярные языки образуют самый нижний уровень иерархии Хомского и являются простейшим классом формальных языков.
Несмотря на свою простоту, это, пожалуй, самый применимый на практике класс: регулярные выражения встроены практически во все языки программирования и инструменты обработки текста.
Исторически понятие конечного автомата восходит к модели нейронных сетей МакКаллока и Питтса (McCulloch, Pitts) 1943 года.
Регулярные выражения были введены Клини (Kleene)~\cite{Kleene56} для описания событий, распознаваемых такими сетями, а строгая теория конечных автоматов и их детерминизация~--- заслуга Рабина и Скотта (Rabin, Scott)~\cite{10.1147/rd.32.0114}.

В этой главе мы последовательно введём три основных способа задания регулярных языков и докажем их эквивалентность.
Сначала будут рассмотрены \emph{регулярные выражения} и \emph{конечные автоматы}~--- как недетерминированные, так и детерминированные.
Затем мы определим \emph{производные Бжозовского} для регулярных выражений, которые будем использовать как инструмент, для построения автомата по выражению.
Далее будут даны взаимные преобразования: \emph{построение конечного автомата по регулярному выражению} и \emph{построение регулярного выражения по конечному автомату}.
Третьим эквивалентным формализмом выступают \emph{лево(право)линейные грамматики}.
Таким образом, в итоге будет установлена цепочка эквивалентности: регулярные выражения $\leftrightarrow$ конечные автоматы $\leftrightarrow грамматики.
Наконец, мы обсудим два фундаментальных инструмента для работы с регулярными языками: \emph{лемму о накачке} и \emph{свойства замкнутости} относительно различных операций.

Регулярные языки находят широчайшее применение на практике: поиск по шаблону в тексте (утилита grep, регулярные выражения в языках программирования%
\sidenote{
    Замечание для программистов.
    Важно понимать, что речь далее пойдёт о формальной конструкции, а не о том, что называется регулярными выражениями в различных языках программирования, библиотеках, инструментах, где под названием \enquote{регулярные выражения} могут скрываться конструкции, существенно более выразительные, чем обсуждаемые здесь, но, тем не менее, родственные.
}), лексический анализ языков программирования, верификация программного обеспечения (model checking), задание шаблонов путей при анализе графов (regular path queries, RPQ) и многие другие области.

Стоит отметить, что рассматриваться будут лишь базовые конструкции, алгоритмы и их свойства.
Фундаментальные результаты можно найти в таких классических учебниках, как <<Introduction to Formal Language Theory>>~\cite{10.5555/578595} и <<Introduction to Automata Theory, Languages, and Computation>>~\cite{Hopcroft+Ullman/79/Introduction}, а для более глубокого изучения области можно обратиться к <<Handbook of Automata Theory>>~\cite{DBLP:books/ems/21/P2021}.


\section{Регулярные выражения}

Регулярные выражения~--- один из часто применяющихся на практике классических способов задать регулярный язык%
\sidenote{
    Замечание для программистов.
    Важно понимать, что речь идёт о формальной конструкции, а не о том, что называется регулярными выражениями в различных языках программирования, библиотеках, инструментах, где под названием \enquote{регулярные выражения} могут скрываться конструкции, существенно более выразительные, чем обсуждаемые здесь.
}.
Основывается этот способ на предложении синтаксиса для описания \emph{регулярных множеств}%
\sidenote{Помним, что язык~--- это множество слов.}.

\begin{definition}[Регулярное множество]
    Регулярное множество (над алфавитом $\Sigma$) это:
    \begin{itemize}
        \item $\varnothing$ или пустое множество;
        \item $\{\varepsilon\}$ или множество, состоящее только из пустой цепочки;
        \item $\{t\}$, $t \in \Sigma$ или множество, состоящее из одного символа из алфавита;
        \item $R_1 \cup R_2$ или объединение множеств, где $R_1$ и $R_2$~--- регулярные множества;
        \item $R_1 \cdot R_2$ или конкатенация множеств, где $R_1$ и $R_2$~--- регулярные множества;
        \item $R^*$ или замыкание Клини, где $R$~--- регулярное множество.
    \end{itemize}
\end{definition}

Для того, чтобы описывать такие множества, можно пользоваться \emph{регулярными выражениями}, которые, фактически, предоставляют синтаксис для описания операций над регулярными множествами, представленными выше.

\begin{definition}[Регулярное выражение]
    Регулярное выражение (над алфавитом $\Sigma$) это:
    \begin{itemize}
        \item пустое множество $\varnothing$;
        \item множество из пустой цепочки $\varepsilon$;
        \item множество из одного символа $t$, $t \in \Sigma$;
        \item объединение регулярных выражений (множеств) $R_1 \mid R_2$, где $R_1$ и $R_2$~--- регулярные выражения;
        \item конкатенация регулярных выражений (множеств) $R_1 \cdot R_2$, где $R_1$ и $R_2$~--- регулярные выражения;
        \item замыкание клини $R^*$, где $R$~--- регулярное выражение;
        \item группирующие скобки $(R)$, где $R$~--- регулярное выражение.
    \end{itemize}
\end{definition}

Отметим несколько важных с прикладной точки зрения моментов.
Во-первых, часто используется расширенный синтаксис, в который добавляются конструкции не увеличивающие выразительную силу, но упрощающие запись.
Например, встречаются следующие расширения%
\sidenote{
    Существуют и другие, однако их мы не будем использовать и, соответственно, рассматривать.
    Читатель может вспомнить, что называется регулярными выражениями в его любимом языке программирования и попробовать самостоятельно выразить имеющиеся там конструкции через базовые.}:
\begin{itemize}
    \item $R? = R \mid \varepsilon$, где $R$~--- регулярное выражение;
    \item $R^+ = R \cdot R^*$, где $R$~--- регулярное выражение.
\end{itemize}

Во-вторых, конструкции $\varnothing$ и $\varepsilon$ используются крайне редко, особенно в случае расширенного синтаксиса, так как часто выражение, эквивалентное использующему данные конструкции, более компактно записывается с использованием расширенного синтаксиса.
В-третьих, оператор конкатенации часто опускается%
\sidenote{Как и знак умножения во многих математических записях.}.

Рассмотрим несколько примеров регулярных выражений.
\begin{example}
    Регулярное выражение $a$ задаёт регулярное множество $\{a\}$ и, соответственно, язык из единственного слова $a$.
\end{example}


\begin{example}
    Регулярное выражение $ab$ задаёт регулярное множество $\{ab\}$ и, соответственно, язык из единственного слова $ab$.
\end{example}


\begin{example}\label{exmpl:regexp_a_star}
    Регулярное выражение $a^*$ задаёт регулярное множество \[R = \bigcup_{i=0}^{\infty}{a^i} = \{\varepsilon, a, aa, aaa, \ldots \}\] и, соответственно, бесконечный язык, содержащий для любого неотрицательного целого $n$ цепочку из символов $a$ длины $n$.
\end{example}

\begin{example}
    Регулярное выражение $a^*b$ задаёт бесконечный язык, содержащий цепочки, полученные конкатенацией цепочек из языка, описанного в примере~\ref{exmpl:regexp_a_star} и символа $b$.
    Иными словами, любая цепочка из языка выглядит как последовательность символов $a$ произвольной длины (возможно пустая), за которой следует символ $b$.
    Примеры цепочек из языка: $b, ab, aaab, aaaaab$.
\end{example}

\begin{example}
    Регулярное выражение $(a\mid b)^*$ задаёт бесконечный язык, цепочки которого имеют произвольную длину и состоят из символов $a$ и $b$ в произвольном порядке.
    Примеры цепочек из языка: $\varepsilon, a, aa, b, ba, aba$.
\end{example}

\begin{example}
    Регулярное выражение\sidenote{Здесь мы воспользовались расширенным синтаксисом.} $(ab)^*c?$ задаёт бесконечный язык, цепочки которого состоят их произвольного количества повторений подпоследовательности $ab$ за которой может следовать (но не обязательно следует) символ $c$.
    Примеры цепочек из языка: $\varepsilon, abab, abc, abababc$.
\end{example}

\section{Конечные автоматы}
	\tikzsetfigurename{RegLang_FA_}

\emph{Конечный автомат} (Finite State Machine, FSM)~--- вычислительная машина, которая имеет конечный набор состояний и может совершать переходы между ними, основываясь на прочитанных входных данных.
Важно отметить, что никакой дополнительной памяти классический конечный автомат не имеет и дополнительных действий (кроме чтения входной ленты) не производит%
\sidenote{Cуществуют автоматы с записью на отдельную (выходную) ленту (например, автоматы Мили~\cite{6771467}) и другие, основанные на конечных автоматах, вычислители, производящие дополнительные действия.
Например, расширенные конечные автоматы (Extended Finite State Machine), которые умеют работать с памятью (переменными)~\cite{Alagar2011, foster.ea:efsm:2020}.}.

Автоматом самого общего вида является \emph{недетерминированный конечный автомат}.
\begin{definition}[Недетерминированный конечный автомат]
    \label{def:NondeterminicticFiniteAutomata}
    \emph{Недетерминированный конечный автомат} (НКА, Nondeterministic Finite Automaton, NFA)~--- это пятёрка $M = \langle Q, Q_S, Q_F, \delta, \Sigma \rangle$, где
    \begin{itemize}
        \item $Q$~--- конечное множество состояний, $Q \neq \varnothing$;
        \item $Q_S \subseteq Q$~--- множество стартовых состояний, $Q_S \neq \varnothing$;
        \item $Q_F \subseteq Q$~--- множество финальных состояний\sidenote{Обратите внимание, множества стартовых и финальных состояний могут пересекаться.}, $Q_F \neq \varnothing$;
        \item $\delta \subseteq Q \times (\Sigma \cup \varepsilon) \times 2^Q$~--- функция переходов, а $\varepsilon \notin \Sigma$;
        \item $\Sigma$~--- конечный алфавит.
    \end{itemize}
\end{definition}

Для поиска путей в графе иногда требуется учитывать не только переходы по исходным меткам рёбер, но и переходы по рёбрам в обратном направлении.
Для этого удобно явно расширить алфавит обратными символами.
Пусть для каждого $a \in \Sigma$ задан символ $a^{-} \notin \Sigma$, причём $(a^{-})^{-} = a$.
Обозначим $\Sigma^{-} = \{a^{-} \mid a \in \Sigma\}$ и $\Sigma^{\leftrightarrow} = \Sigma \cup \Sigma^{-}$.

\mytodo{Дать более аккуратное определение 2-NFA из специализированной профильной литературы.}
\begin{definition}[Двусторонний недетерминированный конечный автомат]
    \label{def:TwoWayNondeterministicFiniteAutomata}
    \emph{Двусторонний недетерминированный конечный автомат} (Two-way Nondeterministic Finite Automaton, 2-NFA)~--- это недетерминированный конечный автомат $M = \langle Q, Q_S, Q_F, \delta, \Sigma^{\leftrightarrow} \rangle$, переходы которого помечены символами из расширенного алфавита $\Sigma^{\leftrightarrow}$.
    Переход по символу $a \in \Sigma$ соответствует чтению метки $a$ в прямом направлении, а переход по символу $a^{-} \in \Sigma^{-}$~--- чтению метки $a$ при движении по ребру графа в обратном направлении.
\end{definition}

Заметим, что функцию переходов можно представить разными способами: это может быть множество троек вида $(q_i, t, q_j)$, матрица, или же граф с метками на рёбрах
\[
G = \langle Q, \{(q_i,t,q_j \mid q_j \in \delta(q_i,t))\}, \Sigma \rangle.
\]
В дальнейшем мы чаще всего будем использовать представление автомата в виде графа.
Чтобы такое представление было полным, в графе отдельно обозначают стартовые и финальные состояния, как это показано на рисунке~\ref{fig:nfa_example}.

Так как нас интересуют конечные автоматы в контексте языков, то будем говорить, что на ленте автомата записано какое-то слово (или строка).
Иными словами, будем говорить, что автомат принимает на вход слово или строку.
При таком подходе автомат естественно рассматривать как \emph{распознаватель}\sidenote{Конечно, на автомат можно смотреть и как на генератор: любой путь от стартового состояния до финального задаёт слово из языка.}.

\begin{example}\label{exmpl:NFA}
    Пусть задан автомат $M=\langle Q, Q_S, Q_F, \delta, \Sigma\rangle$, где
    \begin{itemize}
        \item $Q = \{0,1,2\}$;
        \item $Q_S = \{0\}$;
        \item $Q_F = \{1\}$;
        \item $\Sigma = \{a,b\}$;
        \item $\delta$ определена следующим образом\sidenote{Часто указывают только те случаи, в которых $\delta$ возвращет непустое множество. Для неуказанных входных аргументов считается, что $\delta$ возвращает пустое множество.}:
        \begin{alignat*}{6}
            \delta(0,a) =& \{0,1\}      \qquad    & \delta(1,a) =& \varnothing \qquad & \delta(2,a) =& \varnothing \\
            \delta(0,b) =& \varnothing           & \delta(1,b) =& \{2\}             & \delta(2,b) =& \{1\} \\
            \delta(0,\varepsilon) =& \varnothing & \delta(1,\varepsilon) =& \{0\}   & \delta(2,\varepsilon) =& \varnothing.
        \end{alignat*}
    \end{itemize}
    Тогда его можно представить в виде графа как показано на рисунке~\ref{fig:nfa_example}.
    \begin{marginfigure}
        \begin{center}
            \input{part_02_Foundations/chapter_05_RegularLanguages/figures/02_FiniteAutomata/fig_nfa_example.tex}
        \end{center}
        \caption{Пример конечного автомата в котором состояние $0$~--- стартовое, а состояние $1$~--- финальное}
        \label{fig:nfa_example}
    \end{marginfigure}
\end{example}

Процесс вычислений, проделываемых конечным автоматом, удобно описывать в терминах переходов между \emph{конфигурациями}.

\begin{definition}[Конфигурация]
    Конфигурация $c$ конечного автомата $M = \langle Q, Q_S, Q_F, \delta, \Sigma \rangle$~--- это пара $(q, w)$, где $q\in Q$~--- это текущее состояние автомата, а $w \in \Sigma^*$~--- непросмотренная часть входной строки.
\end{definition}

\begin{definition}[Переход в НКА]
    Будем говорить, что автомат $M = \langle Q, Q_S, Q_F, \delta, \Sigma \rangle$ может перейти из конфигурации $c_1 = (q_1, w_1)$ в конфигурацию $c_2 = (q_2, w_2)$, если
    \[
    c_2 \in \{(q_2,w_2) \mid w_1 = aw_2, (q_1,a, q_2) \in \delta\} \cup \{(q_2,w_1) \mid (q_1, \varepsilon, q_2) \in \delta\}.
    \]
    Обозначать этот факт будем как $c_1 \to c_2$.
    Также будем считать, что на множестве конфигураций задано отношение перехода $(\to):(Q \times \Sigma^*)\times(Q \times \Sigma^*)$.

\end{definition}

\begin{definition}
    Транзитивное рефлексивное замыкание отношения перехода на конфигурациях будем обозначать следующим образом: \[ c_1 \to^* c_2. \]
    Альтернативно, в случае если $c_1 \to^* c_2$, будем говорить, что конфигурация $c_2$ \textit{достижима} из конфигурации $c_1$.
\end{definition}

Для удобства работы с недетерминированными автоматами расширим это отношение на множество конфигураций.

\begin{definition}
Будем говорить, что автомат может перейти из множества конфигураций $C_1$ в множество конфигураций $C_2$, если
\[
C_2 = \bigcup_{c_1 \in C_1} \{c_2 \mid c_1 \to c_2 \}.
\]

Обозначать этот факт будем как  $C_1 \Rightarrow C_2 $.
\end{definition}

\begin{definition}
Транзитивное рефлексивное замыкание отношения перехода на множествах конфигураций будем обозначать следующим образом: \[ C_1 \Rightarrow^* C_2. \]
\end{definition}

Для описания работы автомата $M = \langle Q, Q_S, Q_F, \delta, \Sigma \rangle$ нам понадобятся следующие выделенные типы конфигураций.
\begin{itemize}
    \item Стартовая конфигурация $c_s = (q_s,w)$, $q_s \in Q_S$, $w$~--- цепочка, которая подаётся на вход автомату.
    Для недетерминированного автомата естественно задать множество стартовых конфигураций $C_S = \bigcup_{q_s \in Q_S} (q_s,w)$.
    \item Финальная (принимающая) конфигурация $c_f = (q_f,\varepsilon)$, $q_f \in Q_F$.
\end{itemize}

Таким образом, работу автомата можно описать как последовательность переходов между множествами конфигураций.
Работа начинается с множества стартовых конфигураций и завершается в следующих двух случаях.
\begin{enumerate}
    \item Очередное множество конфигураций содержит финальную конфигурацию:
    \[C_S \Rightarrow^* C_i, c_f \in C_i.\] В этом случае говорят, что автомат \textit{принимает} входную строку.
    \item Очередное множество конфигураций пусто:
    \[C_0 = C_S \Rightarrow^* C_i \Rightarrow^* \varnothing, \text{ для любого } i: c_f \notin C_i.\]
    В этом случае говорят, что автомат \textit{не принимает} или \textit{отвергает} входную строку.
\end{enumerate}

\begin{definition}
    Язык задаваемый недетерминированным конечным автоматом $M = \langle Q, Q_S, Q_F, \delta, \Sigma \rangle$~--- это множество слов, принимаемых им:
     \[ \{w \mid (q_s,w) \to^* (q_f,\varepsilon), q_s \in Q_S, q_f \in Q_F \}.\]
\end{definition}

Так как конфигурация полностью описывает состояние процесса вычислений, то не надо обрабатывать одну и ту же конфигурацию несколько раз.
Это поможет при написании реального интерпретатора: будем отслеживать уже посещённые (обработанные) конфигурации\sidenote{Техника, аналогичная той, что применяется в обходах графов (обход в ширину, обход в глубину) для того, чтобы избежать повторного посещения вершин и, как следствие, зацикливания обхода. Более того, она типична для алгоритмов с рабочим множеством.}.

\begin{example}
    Рассмотрим пример работы конечного автомата, представленного на рисунке~\ref{fig:nfa_example}, на входной цепочке \texttt{abb}.
    В данном автомате стартовое состояние одно, потому множество стартовых конфигураций состоит из одной конфигурации:
    \[
    C_S = \{(0,\texttt{abb})\}.
    \]
    Из состояния \circled{0} существуют два перехода по символу \texttt{a}, который является первым в рассматриваемой цепочке.
    Потому будет получено следующее множество конфигураций:
    \[
        \{(0,\texttt{abb})\} \Rightarrow \{ (0,\texttt{bb}), (1,\texttt{bb})\}.
    \]
    В множестве появилась конфигурация, содержащая финальное состояние.
    Однако цепочка не прочитана до конца, потому продолжаем работу.
    Из нулевого состояния нет переходов по символу \texttt{b}, потому соответствующая конфигурация не породит новых.
    Из второй же конфигурации возможен переход по символу \texttt{b} и по $\varepsilon$:
    \[
        \{ (0,\texttt{bb}), (1,\texttt{bb})\} \Rightarrow \{ (0,\texttt{bb}), (2,\texttt{b})\}
    \]
    И снова, возможен только переход из состояния \circled{2}:
    \[
        \{ (0,\texttt{bb}), (2,\texttt{b})\} \Rightarrow \{(1,\varepsilon)\}.
    \]
    Теперь получена финальная конфигурация: состояние финальное и цепочка полностью обработана.
    Значит можно завершать работу и сообщать, что цепочка принимается автоматом.

    Полностью последовательность переходов выглядит следующим образом:
    \[
        \{(0,\texttt{abb})\} \Rightarrow \{ (0,\texttt{bb}), (1,\texttt{bb})\} \Rightarrow \{ (0,\texttt{bb}), (2,\texttt{b})\} \Rightarrow \{(1,\varepsilon)\}.
    \]
\end{example}

\begin{definition}[Детерминированный конечный автомат]
    \label{def:DeterminicticFiniteAutomata}
    \emph{Детерминированный конечный автомат} (ДКА, Deterministic Finite Automaton, DFA)~--- это пятёрка $M = \langle Q, q_S, Q_F, \delta, \Sigma \rangle$, где
    \begin{itemize}
        \item $Q$~--- конечное множество состояний, $Q \neq \varnothing$;
        \item $q_S \in Q$~--- стартовое состояние;
        \item $Q_F \subseteq Q$~--- множество финальных состояний\sidenote{Обратите внимание, стартовое состояние может входить в множество финальных.}\vspace{1.2\baselineskip}, $Q_F \neq \varnothing$;
        \item $\delta \subseteq Q \times \Sigma \times Q$~--- функция переходов\sidenote{При данном определении автомат будет ещё и \emph{полным}: для каждого состояния и каждого символа определён переход. Часто встречается вариант, когда функция переходов частично определена: для некоторых состояний не определены переходы для некоторых символов.};
        \item $\Sigma$~--- конечный алфавит.
    \end{itemize}
\end{definition}

Основных отличий от недетерминированного автомата два.
Во-первых, стартовое состояние ровно одно.
Во-вторых, функция переходов не допускает переходов по $\varepsilon$ и из любого состояния существует не более одного перехода по конкретному символу.

\begin{example}\label{exmpl:DFA}
    Пусть задан детерминированный автомат $M=\langle Q, q_S, Q_F, \delta, \Sigma\rangle$, где
    \begin{itemize}
        \item $Q = \{0,1,2\}$;
        \item $q_S = 0$;
        \item $Q_F = \{1\}$;
        \item $\Sigma = \{a,b\}$;
        \item $\delta$ определена следующим образом\sidenote{В нашем случае автомат не будет полным, так что $\delta$~--- частично определённая функция.}:
        \begin{alignat*}{6}
            \delta(0,a) =& 1      \qquad    & \delta(1,a) =& 1 \\
            \delta(1,b) =& 2      \qquad    & \delta(2,b) =& 1
        \end{alignat*}
    \end{itemize}
    Тогда его можно представить в виде графа, показанного на рисунке~\ref{fig:dfa_example}.
    \begin{marginfigure}
        \begin{center}
                        \begin{tikzpicture}
                \node[state,initial] (q_0) {$0$};
                \node[state,accepting] (q_1) [right = of q_0] {$1$};
                \node[state] (q_2) [above = of q_1] {$2$};
                \path[->]
                    (q_0) edge[bend left, above]   node {$a$} (q_1)
                    (q_1) edge[loop below, below]   node {$a$} (q_1)
                    (q_1) edge[bend left, left]   node {$b$} (q_2)
                    (q_2) edge[bend left, right]   node {$b$} (q_1);
            \end{tikzpicture}

        \end{center}
        \caption{Пример детерминированного конечного автомата в котором состояние $0$~--- стартовое, а состояние $1$~--- финальное, задающего тот же язык, что и автомат из примера~\ref{exmpl:NFA}}
        \label{fig:dfa_example}
    \end{marginfigure}

\end{example}

\begin{example}
    Рассмотрим пример работы конечного автомата, представленного на рисунке~\ref{fig:dfa_example}, на входной цепочке \texttt{abb}.
    Стартовая конфигурация:
    \[
    (0, \texttt{abb}).
    \]
    Из неё есть переход по символу \texttt{a}.
    Выполним его и получим следующую конфигурацию:
    \[
    (0, \texttt{abb}) \to (1, \texttt{bb}).
    \]
    Наблюдаем конфигурацию с финальным состоянием, но цепочка прочитана не до конца.
    Продолжаем работу:
    \[
    (1, \texttt{bb}) \to (2, \texttt{b}) \to (1, \varepsilon).
    \]
    Теперь конфигурация финальная: цепочка прочитана полностью и мы находимся в финально состоянии.
    Можно сообщить, что автомат принимает эту цепочку.
\end{example}

\section{Производные для регулярных языков}

Производные для регулярных языков предложены в Янушем Бжозовским в работе <<Derivatives of Regular Expressions>> ~\sidecite{Brzozowski1964}.
Будучи достаточно удобным механизмом для решения ряда (в том числе прикладных) задач, производные активно исследовались и мы будем использовать работу~\sidecite{OWENS_REPPY_TURON_2009} как основу для нашего изложения.

Общее определение производных было дано выше и работает, разумеется, и для регулярных языков.
Однако общее определение ничего не говорит о том, как конструктивно вычислить производную.
Для регулярных языков возможно представить алгоритм, который вычисляет производные на основе регулярных выражений.

Прежде чем определить непосредственно производные, определим вспомогательную функцию, которая будет по регулярному выражению определять, принадлежит ли пустая цепочка языку, задаваемому им.
Часто эта функция называется \emph{IsNull}.

\begin{itemize}
    \item $\emph{IsNull}(\varepsilon) = \emph{true}$
    \item $\emph{IsNull}(\varnothing) = \emph{false}$
    \item $\emph{IsNull}(x) = \emph{false}$
    \item $\emph{IsNull}(R_1 \cdot R_2) = \emph{IsNull}(R_1) \ \& \ \emph{IsNull}(R_2) $
    \item $\emph{IsNull}(R_1 \mid R_2) = \emph{IsNull}(R_1) \mid \emph{IsNull}(R_2) $
    \item $\emph{IsNull}(R^*) = \emph{true}$
\end{itemize}

Далее, введём функцию $\mathcal{V}$:
\[
    \mathcal{V}(R) =
    \begin{cases}
        \varepsilon, \emph{IsNull}(R) \\
        \varnothing, \text{ иначе}
    \end{cases}.
\]

Теперь мы готовы определить производные для различных конструкций регулярных выражений.
\begin{itemize}
    \item $\partial_t(\varepsilon) = \varnothing$
    \item $\partial_t(\varnothing) = \varnothing$
    \item $\partial_t(t) = \varepsilon$
    \item $\partial_t(x) = \varnothing \text{ для } t \neq x$
    \item $\partial_t(R_1 \cdot R_2) = \partial_t(R_1) \cdot (R_2) \mid \mathcal{V}(R_1) \cdot \partial_t(R_2)$
    \item $\partial_t(R_1 \mid R_2) = \partial_t(R_1) \mid \partial_t(R_2)$
    \item $\partial_t(R^*) =$\sidenote{Интересное упражнение~--- показать это, расписав по определению звезду Клини.} $ \partial_t(R)\cdot R^*$.
\end{itemize}

\begin{example}
    Вычислим производную $a^*(a \mid b)$ по символу $a$.
    \[
        \partial_a(a^*(a \mid b)) = \partial_a(a^*) (a \mid b) \mid \mathcal{V}(a^*) \cdot \partial_a(a \mid b) \\
    \]
    Вычислим отдельные подвыражения.
    \begin{align*}
        \partial_a(a^*) & = \partial_a(a) a^* = a^* \\
        \partial_a(a \mid b) & = \partial_a(a) \mid \partial_a(b) = \varepsilon \mid \varnothing = \varepsilon
    \end{align*}
    Соединим всё вместе.
    \begin{align*}
        \partial_a(a^*(a \mid b)) & = \partial_a(a^*) (a \mid b) \mid \mathcal{V}(a^*) \cdot \partial_a(a \mid b) \\
                                  & = a^* (a \mid b) \mid \varepsilon \cdot \varepsilon \\
                                  & = a^* (a \mid b) \mid \varepsilon \\
                                  & = (a^* (a \mid b))?
    \end{align*}
\end{example}

Одной из прикладных задач, которую можно решить с помощью производных, является принадлежность строки языку, задаваемому регулярным выражением\sidenote{Что важно, без преобразований регулярного выражения в конечный автомат, как было показано выше.}.
Пусть дано регулярное выражение $R$ и нужно проверить, принадлежит ли строка $w=t_0 \ldots \ t_{n_1}$ языку, задаваемому этим выражением.
Для этого необходимо вычислить следующее выражение:
\[
\emph{IsNull}(\partial_{t_{n-1}}(\ldots(\partial_{t_0}(R))\ldots)).
\]
Если его значение \emph{true}, то строка принадлежит языку.
Иначе не принадлежит.

\begin{example}
    Пусть дано регулярное выражение $a(a \mid (b \ b))^*$ и строка \texttt{abb}.
    Принадлежит ли она языку, задаваемому данным выражением?
    Для ответа на этот вопрос вычислим следующее выражение:
    \[
        \emph{IsNull}(\partial_b(\partial_b(\partial_a(a^*(a \mid (b \ b))^*)))).
    \]
    Для начала вычислим производные.
    \begin{align*}
        \partial_a(a^*(a \mid (b \ b))^*) &= \partial_a(a) (a \mid (b \ b))^* \mid \mathcal{V}(a)(\partial_a(a \mid (b \ b))^*) \\
                                          &= \varepsilon (a \mid (b \ b))^* \mid \varnothing (\partial_a(a \mid (b \ b))^*) \\
                                          &= (a \mid (b \ b))^* \mid \varnothing \\
                                          &= (a \mid (b \ b))^*\\
        \partial_b((a \mid (b \ b))^*)    &= \partial_b((a \mid (b \ b))) (a \mid (b \ b))^* \\
                                          &= (\varnothing \mid b) (a \mid (b \ b))^* \\
                                          &= b (a \mid (b \ b))^* \\
        \partial_b(b (a \mid (b \ b))^*)  &= (a \mid (b \ b))^*
    \end{align*}
    $\emph{IsNull}((a \mid (b \ b))^*) = \emph{true}$ по определению.
    Таким образом, цепочка \texttt{abb} принадлежит языку, задаваемому регулярным выражением $a(a \mid (b \ b))^*$.
\end{example}

\section{Построение конечного автомата по регулярному выражению}
\label{sec:RegexToFA}
\tikzsetfigurename{RegLang_RegexToFA_}

Конечные автоматы и регулярные выражения~--- два способа задать один и тот же класс языков.
Чтобы убедиться в их равной выразительной силе, начнём с построения конечного автомата по регулярному выражению.
Использование производных позволит нам сразу получить \emph{полный детерминированный автомат}%
\sidenote{Существуют и другие алгоритмы построения автомата по регулярному выражению, например, алгоритм Томпсона~\cite{10.1145/363347.363387} или алгоритм Глушкова~\cite{Glushkov1961}.
Каждый из них имеет свои особенности: некоторые строят недетерминированные автоматы, другие используют $\varepsilon$-переходы.}.

Состояниям автомата будем сопоставлять регулярные выражения.
Если выбрать некоторое состояние автомата в качестве стартового, то получившийся автомат будет распознавать язык, задаваемый соответствующим регулярным выражением.

Чтобы избежать дублирования состояний, необходимо проверять эквивалентность языков, которые задаются этими выражениями.
Хотя задача проверки эквивалентности регулярных языков разрешима, на практике в данном алгоритме используют более простую \emph{структурную эквивалентность} регулярных выражений%
\sidenote{При проверке структурной эквивалентности обычно учитывают ассоциативность, коммутативность и идемпотентность некоторых операций, чтобы избежать избыточных различий.}~\sidecite{OWENS_REPPY_TURON_2009}.


Таким образом, можно сформулировать следующий алгоритм построения автомата по регулярному выражению $R$.

\begin{enumerate}
    \item В качестве стартового состояния берём выражение $R$ и добавляем его в рабочее множество $W$.
    \item Пока $W$ не пусто:
    \begin{enumerate}
        \item Извлекаем очередное состояние $q$, которому соответствует регулярное выражение $R_q$.
        \item Для каждого символа $a \in \Sigma$ вычисляем производную $\partial_a(R_q)$.
        \item Если полученное выражение задаёт новый (ранее не встречавшийся) язык, добавляем соответствующее состояние.
        \item Если $\emph{IsNull}(R_q) = \emph{true}$, помечаем состояние как финальное.
        \item Добавляем ребро $q \xrightarrow{a} q'$ для каждого вычисленного перехода.
    \end{enumerate}
\end{enumerate}

В результате получим полный детерминированный конечный автомат\sidenote{Формально, если честно проверять равенство языков, можно получить минимальный, но на практике это сложно. Но часто всё же достаточно близко к минимальному получается.}.

\begin{example}
    Рассмотрим построение автомата для регулярного выражения $a(a \mid (b \ b))^*$.

    Стартовому состоянию соответствует выражение $a(a \mid (b \ b))^*$.
    Так как $\emph{IsNull}(a(a \mid (b \ b))^*) = \emph{false}$, это состояние не является финальным.
    Вычислим производные по всем символам из алфавита:
    \begin{align*}
        \partial_a(a(a \mid (b \ b))^*) & = (a \mid (b \ b))^*, \\
        \partial_b(a(a \mid (b \ b))^*) & = \varnothing.
    \end{align*}
    В результате появляются два новых состояния: $(a \mid (b \ b))^*$ и $\varnothing$.
    Второе состояние соответствует пустому языку и является \emph{поглощающим} (или \emph{дьявольским}) состоянием, а первое~--- финальным, так как $\emph{IsNull}((a \mid (b \ b))^*) = \emph{true}$.
    На рисунке~\ref{fig:regexp_to_dfa_example_step_1} показан автомат после первого шага алгоритма.

    \begin{marginfigure}
        \begin{center}
            \scalebox{0.64}{
                        \begin{tikzpicture}
                \node[elliptic state,initial] (q_0) {$a(a \mid (b \ b))^*$};
                \node[elliptic state,accepting] (q_1) [right = of q_0] {$(a \mid (b \ b))^*$};
                \node[state] (q_2) [above = of q_0] {$\varnothing$};
                \path[->]
                    (q_0) edge[bend right, below]   node {$a$} (q_1)
                    (q_0) edge[bend left, left]   node {$b$} (q_2);
            \end{tikzpicture}

            }
        \end{center}
        \caption{Построение автомата по регулярному выражению с помощью производных: первый шаг}
        \label{fig:regexp_to_dfa_example_step_1}
    \end{marginfigure}

    На следующем шаге вычислим производные для состояния $(a \mid (b \ b))^*$:
    \begin{align*}
        \partial_a((a \mid (b \ b))^*) & = (a \mid (b \ b))^*, \\
        \partial_b((a \mid (b \ b))^*) & = b(a \mid (b \ b))^*.
    \end{align*}
    Таким образом, появляется новое состояние $b(a \mid (b \ b))^*$, а автомат принимает вид, показанный на рисунке~\ref{fig:regexp_to_dfa_example_step_2}.

    \begin{marginfigure}
        \begin{center}
            \scalebox{0.64}{
            \input{part_02_Foundations/chapter_05_RegularLanguages/figures/04_RegexToFA/fig_regexp_to_dfa_example_step_2.tex}
            }
        \end{center}
        \caption{Построение автомата по регулярному выражению с помощью производных: второй шаг}
        \label{fig:regexp_to_dfa_example_step_2}
    \end{marginfigure}

    Осталось обработать состояние $b(a \mid (b \ b))^*$:
    \begin{align*}
        \partial_a(b(a \mid (b \ b))^*) & = \varnothing, \\
        \partial_b(b(a \mid (b \ b))^*) & = (a \mid (b \ b))^*.
    \end{align*}
    В результате добавляются новые переходы, и автомат принимает окончательный вид (рисунок~\ref{fig:regexp_to_dfa_example_step_3}).

    \begin{marginfigure}
        \begin{center}
            \scalebox{0.64}{
                        \begin{tikzpicture}
                \node[elliptic state,initial] (q_0) {$a(a \mid (b \ b))^*$};
                \node[elliptic state,accepting] (q_1) [right = of q_0] {$(a \mid (b \ b))^*$};
                \node[state] (q_2) [above = of q_0] {$\varnothing$};
                \node[elliptic state] (q_3) [above = of q_1] {$b(a \mid (b \ b))^*$};
                \path[->]
                    (q_0) edge[bend right, below]   node {$a$} (q_1)
                    (q_0) edge[bend left, left]   node {$b$} (q_2)
                    (q_1) edge[bend left, left]   node {$b$} (q_3)
                    (q_2) edge[loop left, left]   node {$b,a$} (q_2)
                    (q_1) edge[loop below, above]   node {$a$} (q_1)
                    (q_3) edge[bend left, right]   node {$b$} (q_1)
                    (q_3) edge[bend left, below]   node {$a$} (q_2);
            \end{tikzpicture}

            }
        \end{center}
        \caption{Построение автомата по регулярному выражению с помощью производных: результирующий автомат}
        \label{fig:regexp_to_dfa_example_step_3}
    \end{marginfigure}
\end{example}

В результате построения получен автомат, близкий по структуре к автомату из примера~\ref{exmpl:DFA}.
Разница лишь в том, что приведённый здесь автомат является полным.
Само же сходство не случайно: автоматы действительно задают один и тот же регулярный язык.

\section{Построение регулярного выражения по конечному автомату}

Для построения регулярного выражения по конечному будем использовать вариацию классической конструкции, часто связываемой с теоремой Клини о эквивалентности языков, задаваемых регулярными выражениями и конечными автоматами, для доказательства которой нужно, в частности, построить регулярное выражение по конечному автомату.
Регулярное выражение будем строить по недетерминированному автомату специального вида: потребуем, чтобы у него было ровно одно стартовое состояние и ровно одно финальное%
\sidenote{Любой автомат можно привести к такому виду.
Предположим, что в автомате два стартовых состояния $q_0^s$ и $q_1^s$.
Добавим новое состояние $q_s$, назначим его стартовым, добавим переходы $q_s \xrightarrow{\varepsilon} q_0^s$ и $q_s \xrightarrow{\varepsilon} q_1^s$, уберём $q_0^s$ и $q_1^s$ из стартовых.
Аналогично можно поступить и с финальными состояниями.}.

Будем в цикле выполнять последовательно две операции.
\begin{enumerate}
    \item Объединение параллельных рёбер.
    \item Устранение состояния $v$. За один шаг можем устранить любое состояние кроме стартового или финального.
\end{enumerate}
Цикл повторяется до тех пор, пока в автомате не останется ровно два состояния: стартовое и финальное.
После завершения цикла необходимо ещё раз объединить параллельные рёбра.
После этого, по полученному автомату можно построить итоговое регулярное выражение.

Параллельные рёбра вида $p_i \xrightarrow{R_1} q_i, \ldots, p_i \xrightarrow{R_k} q_i$ объединяются в одно ребро вида $p_i \xrightarrow{R_1 \mid \ldots \mid R_k} q_i$, метка которого~--- объединение всех регулярных выражений с параллельных рёбер.
Пример такого объединения приведён на рисунках~\ref{fa:fa1} и~\ref{fa:fa2}.


\begin{marginfigure}
    \begin{center}
        \begin{tikzpicture}

            \node[state] (q_0)                   {$q_i$};
            \node[state] (q_1) [right = 1.8 of q_0]  {$q_j$};

            \path[->]
                (q_0) edge[bend left = 45, above]   node {$R_1$} (q_1)
                (q_0) edge[above]   node {$R_2$} (q_1)
                (q_0) edge[bend right, below]  node {$R_3$} (q_1);

        \end{tikzpicture}
    \end{center}
    \caption{Два состояния с параллельными рёбрами}
    \label{fa:fa1}
\end{marginfigure}
\begin{marginfigure}

    \begin{center}
        \begin{tikzpicture}

            \node[state] (q_0)                    {$q_i$};
            \node[state] (q_1) [right = 1.8 of q_0]  {$q_j$};

            \path[->]
                (q_0) edge[above]  node {$R_1 \mid R_2 \mid R_3$} (q_1);

        \end{tikzpicture}
    \end{center}
    \caption{Результат объединения параллельных рёбер для состояний, представленных на рисунке~\ref{fa:fa1}}
    \label{fa:fa2}
\end{marginfigure}

Операция устранения состояния $v$ устроена следующим образом.
Рассмотрим три типа рёбер.
\begin{enumerate}
    \item Рёбра вида $p_i \xrightarrow{R_{p_i}} v$, где $p_i \neq v$.
    \item Рёбра вида $v \xrightarrow{R_{q_j}} q_i$, где $q_i \neq v$.
    \item Рёбра вида\sidenote{Рёбер такого вида для заданной вершины всегда не более одного, так как мы производим объеднинение параллельных рёбер на каждой итерации алгоритма.} $v \xrightarrow{R_v} v$.
\end{enumerate}

В результате устранения состояния $v$ все рассмотренные рёбра должны быть удалены.
Взамен них должны быть добавлены рёбра вида $p_i \xrightarrow{R_{p_i} \cdot R_v^* \cdot R_{q_j}} q_j$.
Пример описанного преобразования представлен на рисунках~\ref{fa:fa3} и~\ref{fa:fa4}

\begin{figure}
    \begin{center}
        \begin{subfigure}[b]{0.47\textwidth}
            \begin{center}

                \begin{tikzpicture}

                \node[state] (p_0)          {$p_0$};
                \node[text width=0.3cm]  (p_1) [below = 0.3 of p_0] {$\vdots$};
                \node[state] (p_2) [below = 0.3 of p_1]  {$p_i$};
                \node[text width=0.3cm] (p_3) [below = 0.3 of p_2]  {$\vdots$};
                \node[state] (v_0) [right = of  p_2] {$v$};
                \node[state] (q_2) [right = of v_0] {$q_j$};
                \node[text width=0.3cm] (q_1) [above = 0.3 of  q_2]  {$\vdots$};
                \node[state] (q_0) [above = 0.3 of q_1]  {$q_0$};
                \node[text width=0.3cm] (q_3) [below = 0.3 of q_2]  {$\vdots$};

                \path[->]
                    (p_0) edge[above]  node {$R_{p_0}$} (v_0)
                    (p_2) edge[below]  node {$R_{p_i}$} (v_0)
                    (v_0) edge[left]  node {$R_{q_0}$} (q_0)
                    (v_0) edge[below]  node {$R_{q_j}$} (q_2)
                    (v_0) edge[loop above, above]  node {$R_v$} (v_0);

                \end{tikzpicture}
            \end{center}
            \caption{Автомат перед устранением состояния $v$}
            \label{fa:fa3}

        \end{subfigure}\hspace{5mm}
        \begin{subfigure}[b]{0.47\textwidth}
            \begin{center}

                \begin{tikzpicture}

                    \node[state] (p_0)          {$p_0$};
                    \node[text width=0.3cm]  (p_1) [below = 0.3 of p_0] {$\vdots$};
                    \node[state] (p_2) [below = 0.3 of p_1]  {$p_i$};
                    \node[text width=0.3cm] (p_3) [below = 0.3 of p_2]  {$\vdots$};

                    \node[text width=0.3cm] (v_0) [right = of p_2] {};

                    \node[state] (q_2) [right = of v_0] {$q_j$};
                    \node[text width=0.3cm] (q_1) [above = 0.3 of q_2]  {$\vdots$};
                    \node[state] (q_0) [above = 0.3 of q_1]  {$q_0$};
                    \node[text width=0.3cm] (q_3) [below = 0.3 of q_2]  {$\vdots$};

                    \path[->]
                    (p_0) edge[bend left, above]  node {$R_{p_0} R_v^* R_{q_0}$} (q_0)
                    (p_0) edge[bend left, left]  node {$R_{p_0} R_v^* R_{q_j}$} (q_2)
                    (p_2) edge[bend right, left]  node {$R_{p_i} R_v^* R_{q_0}$} (q_0)
                    (p_2) edge[bend right, below]  node {$R_{p_0} R_v^* R_{q_0}$} (q_2);

                \end{tikzpicture}

            \end{center}

            \caption{Результат устранения состояния $v$ для автомата, представленного на рисунке~\ref{fa:fa3}}
            \label{fa:fa4}
        \end{subfigure}
    \end{center}
    \caption{Общая схема устранения состояния $v$ при построении регулярного выражения по конечному автомату}
    \label{fa:fa_3_4}
\end{figure}

После завершения основного цикла (когда в автомате осталось ровно два состояния), необходимо ещё раз объединить параллельные рёбра.
Общий вид получившегося автомата представлен на рисунке~\ref{fig:fa_to_regexp_final}.
По такому автомату строим выражение вида \[(R_1^* \cdot (R_2 \cdot R_3^* \cdot R_4 \cdot R_1^*)^* \cdot R_2 \cdot R_3^*,\] которое и будет ответом.
\begin{marginfigure}

    \begin{center}
        \begin{tikzpicture}

            \node[state, initial] (q_0)          {$0$};
            \node[state, accepting] (q_1) [right = of q_0]  {$1$};
            \path[->]
                (q_0) edge[bend left, above]  node {$R_2$} (q_1)
                (q_1) edge[bend left, below]  node {$R_4$} (q_0)
                (q_1) edge[loop above, above]  node {$R_3$} (q_1)
                (q_0) edge[loop above, above]  node {$R_1$} (q_0);

        \end{tikzpicture}

    \end{center}

    \caption{Общий вид автомата после завершения основного цикла исключения вершин при построении по нему регулярного выражения}
    \label{fig:fa_to_regexp_final}

\end{marginfigure}

\begin{example}
    Построим регулярное выражение по автомату, представленному на рисунке~\ref{fig:nfa_example}.
    В автомате уже одно финальное и одно стартовое состояние.
    При этом, существует всего одно состояние, которое не является ни финальным, ни стартовым.
    Это состояние \circled{2}.
    Его и устраним.

    После устранения состояния \circled{2} получим автомат, представленный на рисунке~\ref{fig:nfa_to_regexp}.
    Теперь наш автомат принял вид, соответствующий представленному на рисунке~\ref{fig:fa_to_regexp_final}, и можно выписать результирующее регулярное выражение, которое будет иметь следующий вид%
    \sidenote{Конечно, получившееся выражение не является минимальным: достаточно легко можно придумать более простое и компактное регулярное выражение, задающее тот же самый язык.
    Но алгоритм нам и не обещал никаких дополнительных свойств получившегося выражения.}:
    \[
    a^*(a \ (b \ b)^* \ \varepsilon \ a^*)^* \ a \ (b \ b)^* =
    a^*(a \ (b \ b)^* \ a^*)^* \ a \ (b \ b)^*.
    \]

    \begin{marginfigure}
        \begin{center}
            \begin{tikzpicture}
                \node[state,initial] (q_0) {$0$};
                \node[state,accepting] (q_1) [right = of q_0] {$1$};
                \path[->]
                    (q_0) edge[bend left, above]   node {$a$} (q_1)
                    (q_0) edge[loop above, above]   node {$a$} (q_0)
                    (q_1) edge[loop above, above]   node {$bb$} (q_1)
                    (q_1) edge[bend left, below]  node {$\varepsilon$} (q_0);
            \end{tikzpicture}
        \end{center}
        \caption{Пример построения регулярного выражения по конечному автомату: финальное состояние автомата}
        \label{fig:nfa_to_regexp}
    \end{marginfigure}

\end{example}

\section{Лево(право)линейные грамматики}

Наложив некоторые ограничения на внешний вид правил грамматики можно получить грамматики, задающие регулярные языки.

\begin{definition}[Леволинейная грамматика]
    Грамматика $G=\langle \Sigma, N, P, S \rangle$ называется леволинейной, если все её правила имеют вид
    \[N_i \to \alpha w,\]
    где $N_i \in N$, $\alpha \in \{\varepsilon\} \cup N$, $w \in \Sigma ^*$.
\end{definition}

\begin{definition}[Праволинейная грамматика]
    Грамматика $G=\langle \Sigma, N, P, S \rangle$ называется праволинейной, если все её правила имеют вид
    \[N_i \to  w \alpha,\]
    где $N_i \in N$, $\alpha \in \{\varepsilon\} \cup N$, $w \in \Sigma ^*$.
\end{definition}

Ноам Хомский и Джордж Миллер в работе~\sidecite{chomsky1958finite} показали, что лево(право)линейные грамматики задают регулярные языки.
Приведём процедуры построения автомата по грамматике и наоборот, грамматики по автомату.

Пусть дан конечный автомат $M = \langle \Sigma, Q, q_s, Q_f, \delta \rangle$. По нему можно построить праволинейную грамматику $G=\langle \Sigma, N, S, P \rangle$, где
\begin{itemize}
    \item $N = Q$
    \item $P = \{ q_i \to t q_j \mid (q_i, t, q_j)\in \delta\} \cup \{ q_i \to \varepsilon \mid q_i \in Q_F\}$
    \item $S = q_s$
\end{itemize}

Аналогичным образом строится автомат по праволинейной грамматике.
Упростить процедуру можно если заранее привести правила к виду $N_i \to tN_j$, где $t\in \Sigma$, добавив необходимое количество новых нетерминалов:
правило вида $N_i \to twN_k$ преобразуется в два правила
\begin{align*}
    N_i & \to tN_l  \\
    N_l & \to wN_k,
\end{align*}
а если цепочка $w$ содержит более одного символа, то преобразование продолжается аналогичным образом для нетерминала $N_l$ (и далее), пока все правила не примут нужный вид.

\begin{example}[Построение грамматики по автомату]
Построим грамматику по автомату, представленному на рисунке~\ref{fig:dfa_example}.
В автомате три состояния, соответственно, в грамматике будет три нетерминала: $Q_0, Q_1, Q_2$.
Нетерминал $Q_0$ будет стартовым, так как он соответствует стартовому состоянию.
Правила грамматики формируются на основе переходов в автомате и правил для финальных состояний, и в нашем случае будут следующими:
\begin{align*}
    Q_0 \to a Q_1  \quad & \quad Q_1 \to a Q_1 \\
    Q_1 \to b Q_2  \quad & \quad Q_2 \to b Q_1 \\
                     \quad & \quad Q_1 \to \varepsilon.
\end{align*}
Таким образом, мы получили грамматику, которая порождает тот же язык, что и исходный автомат%
\sidenote{В качестве упражнения можно взять несколько цепочек, принадлежащих языку, и проследить, как переходы в автомате в процессе раскознования цепочки соотносятся с применением правил грамматики при её выоде.}.
\end{example}

\begin{marginfigure}
    \begin{center}
                \begin{tikzpicture}
            \node[state,initial] (S) {$S$};
            \node[state,accepting,right=of S] (N1) {$N_1$};
            \node[state,above=of N1] (N2) {$N_2$};

            \path[->]
                (S) edge[bend left, above]   node{$a$} (N1)
                (N1) edge[loop below, below] node{$a$} (N1)
                (N1) edge[bend left, left]   node{$b$} (N2)
                (N2) edge[bend left, right]  node{$b$} (N1);
        \end{tikzpicture}

    \end{center}
    \caption{Автомат, построенный по грамматике $S \to aN_1$, $N_1 \to aN_1 \mid bbN_1 \mid \varepsilon$.
    Эквивалентен автомату из примера~\ref{exmpl:DFA}.}
    \label{fig:grammar_to_dfa}
\end{marginfigure}

\begin{example}[Построение автомата по грамматике]
Рассмотрим грамматику:
\begin{align*}
    S \to a N_1, \quad
    N_1 \to a N_1 \mid b b N_1 \mid \varepsilon.
\end{align*}
Построим по ней конечный автомат.

Первым шагом разобьём продукцию $N_1 \to b b N_1$ на две, введя дополнительный нетерминал.
Новая грамматика будет иметь следующий вид.

\begin{align*}
    S \to a N_1, \quad
    N_1 \to a N_1 \mid b N_2 \mid \varepsilon , \quad
    N_2 \to b N_1.
\end{align*}

Теперь можно заняться конструированием автомата.
Множество состояний автомата соответствует нетерминалам грамматики:
\begin{align*}
    Q = \{S, N_1, N_2\}.
\end{align*}
Алфавит: $\Sigma = \{a, b\}$.
Переходы автомата определяются на основе правил грамматики:
\begin{align*}
\delta = \{
(S, a, N_1),\;
(N_1, a, N_1),\;
(N_1, b, N_2),\;
(N_2, b, N_1)
\}.
\end{align*}
Финальные состояния~--- те, для которых в грамматике есть правило с правой частью $\varepsilon$:
\begin{align*}
Q_F = \{ N_1 \}.
\end{align*}
Таким образом, получаем автомат
\begin{align*}
M = \langle \Sigma, Q, S, Q_F, \delta \rangle,
\end{align*}
распознающий язык всех цепочек, начинающихся с символа $a$ и содержащих произвольное количество $a$ и пар $bb$ после него.

Построенный автомат показан на рисунке~\ref{fig:grammar_to_dfa}. Обратите внимание, что этот автомат эквивалентен автомату из примера~\ref{exmpl:DFA}, только обозначения состояний изменены: $S \leftrightarrow 0$, $N_1 \leftrightarrow 1$, $N_2 \leftrightarrow 2$.
\end{example}

\section{Лемма о накачке}

Лемма о накачке для регулярных языков позволяет проверить, что заданный язык не является регулярным.

\begin{lemma}
    Пусть $L$~--- регулярный язык над алфавитом $\Sigma$, тогда существует такое $n$, что для любого слова $\omega \in L$, $|\omega| \geq n$ найдутся слова $x,y,z\in \Sigma^*$, для которых верно: $xyz = \omega, y\neq \varepsilon,|xy|\leq n$ и для любого $k \geq 0$  $xy^kz \in L$.
\end{lemma}

\begin{proofSketch}
    \begin{enumerate}
        \item Так как язык регулярный, то для него можно построить конечный автомат $M = \langle Q, q_s,Q_f, \delta, \Sigma \rangle$.
              В том числе, минимальный по количеству состояний.
        \item В качестве $n$ возьмём $|Q| + 1$.
        \item Легко заметить, что для любой цепочки $w \in L, |w| > n$ путь в автомате, соответствующий принятию данной цепочки, будет содержать хотя бы один цикл.
              Действительно, в ориентированном графе с $k$ вершинами (а именно таким является автомат по построению) максимальная длина пути без повторных посещений вершин (соответственно, без циклов) не больше $k - 1$.
        \item Выберем любой цикл. Он будет задавать искомые цепочки $x, y$ и $z$ так, как представлено на рисунке~\ref{fig:reg_lang_pumping_lemma}.
              Заметим, что вход в цикл и выход из него в общем случае могут не совпадать, что даёт несколько вариантов разбиения пути на части, и на рисунке представлен лишь один из возможных.
              \qedhere
    \end{enumerate}
\end{proofSketch}

\begin{figure}
    \caption{Иллюстрация идеи доказательства леммы о накачке для регулярных языков: любой путь в графе, длина которого достаточно большая, может быть разбит на три части из леммы ($x$~--- красный подпуть, $y$~--- синий, $z$~--- зелёный), а многократный проход по циклу $y$ позволяет \enquote{накачать} слово.}
    \label{fig:reg_lang_pumping_lemma}
    \begin{center}
        \begin{tikzpicture}[->]
            \node[state, initial] (q1) {$q_1$};
            \node[state, right = 2 of q1] (q2) {$q_2$};
            \node[state, accepting, right = 2 of q2] (q3) {$q_3$};
            \node[state, above = 2 of q2] (q4) {$q_4$};
            \draw (q1) edge[above, snake it] node{} (q2)
            (q2) edge[bend right, snake it, side by side={red}{blue}] node{} (q4)
            (q4) edge[bend right, snake it, color=blue] node[left]{$y$} (q2)
            (q4) edge[above, snake it, color=green] node[right]{$z$} (q3)
            (q1) edge[above, snake it, color=red] node{$x$} (q2)
            ;
        \end{tikzpicture}
    \end{center}

\end{figure}

\section{Замкнутость регулярных языков относительно теоретико-множественных операций}
\label{sec:RegularClosureProperties}
\tikzsetfigurename{RegLang_Closure_}

\begin{theorem}
    Класс регулярных языков над алфавитом $\Sigma$ замкнут относительно следующих теоретико-множественных операций:
    \begin{enumerate}
        \item Объединение: если $L_1, L_2$~--- регулярные, то $L_1 \cup L_2$~--- регулярный;
        \item Пересечение: если $L_1, L_2$~--- регулярные, то $L_1 \cap L_2$~--- регулярный;
        \item Дополнение: если $L$~--- регулярный, то $\overline{L}$~--- регулярный;
        \item Разность: если $L_1, L_2$~--- регулярные, то $L_1 \setminus L_2$~--- регулярный.
    \end{enumerate}
\end{theorem}

\begin{proof}
    Пусть $L_1$ и $L_2$ распознаются конечными автоматами над одним алфавитом $\Sigma$:
    \[
        M_1 = \langle Q^1, Q_S^1, Q_F^1, \delta^1, \Sigma \rangle, \qquad
        M_2 = \langle Q^2, Q_S^2, Q_F^2, \delta^2, \Sigma \rangle.
    \]


    \paragraph{1. Объединение.}
    Построим недетерминированный автомат с $\varepsilon$-переходами,
    вводя новое начальное состояние $q_0$, из которого по $\varepsilon$ можно попасть
    в начальные состояния обоих автоматов:
    \[
        M = \langle Q^1 \cup Q^2 \cup \{q_0\},\; \{q_0\},\; Q_F^1 \cup Q_F^2,\; \delta,\; \Sigma \rangle,
    \]
    где $\delta$ совпадает с $\delta^1$ и $\delta^2$ на соответствующих состояниях, а дополнительно
    \[
        \delta(q_0, \varepsilon) = Q_S^1 \cup Q_S^2.
    \]
    Тогда $L(M) = L_1 \cup L_2$.

    \paragraph{2. Пересечение.}
    Пусть автоматы детерминированы. Построим автомат произведения:
    \[
        M_3 = \langle Q^1 \times Q^2,\; Q_S^1 \times Q_S^2,\; Q_F^1 \times Q_F^2,\; \delta^3,\; \Sigma \rangle,
    \]
    где переходы определяются как:
    \[
        \delta^3\big((q^1, q^2), t\big) = \big(\delta^1(q^1, t),\; \delta^2(q^2, t)\big).
    \]
    Тогда%
    \sidenote{Аккуратно доказать это можно расписав переходы между конфигурациями соответствующих автоматов при обработке цепочки.}
    \[
        w \in L(M_3) \iff w \in L_1 \text{ и } w \in L_2,
    \]
    следовательно, $L(M_3) = L_1 \cap L_2$.

    \paragraph{3. Дополнение.}
    Пусть $M = \langle Q, \{q_S\}, Q_F, \delta, \Sigma \rangle$~--- детерминированный полный автомат.
    Построим автомат
    \[
        M' = \langle Q, \{q_S\}, Q \setminus Q_F, \delta, \Sigma \rangle.
    \]
    Очевидно, $L(M') = \overline{L(M)}$.

    \paragraph{4. Разность.}
    Разность выражается через пересечение и дополнение:
    \[
        L_1 \setminus L_2 = L_1 \cap \overline{L_2}.
    \]
    Регулярные языки замкнуты относительно обеих операций, следовательно, и относительно разности тоже замкнуты.
\end{proof}

%\begin{example}[Конструкция произведения]\label{ex:product-automaton}
%    Пусть $M_1$ и $M_2$~--- детерминированные автоматы над алфавитом $\Sigma=\{a,b\}$.
%    В автомате $M_1$ есть переходы
%    $q^1 \xrightarrow{a} p^1$, $q^1 \xrightarrow{b} p^1$ и $p^1 \xrightarrow{b} p^1$,
%    а в автомате $M_2$~--- переходы
%    $q^2 \xrightarrow{a} q^2$ и $q^2 \xrightarrow{b} p^2$.
%    Тогда в автомате произведения $M_1 \times M_2$ возникают переходы:
%    \begin{align*}
%        (q^1, q^2) \xrightarrow{a} (p^1, q^2), \quad
%        (p^1, q^2) \xrightarrow{b} (p^1, p^2), \quad
%        (q^1, q^2) \xrightarrow{b} (p^1, p^2),
%    \end{align*}
%    что проиллюстрировано на рис.~\ref{fig:product-automaton-example}.
%    Принимающими являются пары из принимающих состояний обоих автоматов.
%\end{example}
%
%
%\begin{marginfigure}
%    \centering
%    \scalebox{0.64}{
%    \begin{tikzpicture}
%        \node[elliptic state,initial] (q0) {$(q^1, q^2)$};
%        \node[elliptic state] (p0) [right=of q0] {$(p^1, q^2)$};
%        \node[elliptic state,accepting] (p1) [below right=of p0] {$(p^1, p^2)$};
%
%        \path[->]
%            (q0) edge[bend left, above] node {$a$} (p0)
%            (p0) edge[bend left, right] node {$b$} (p1)
%            (q0) edge[bend right, below left] node[swap] {$b$} (p1);
%    \end{tikzpicture}}
%    \caption{Пример переходов в произведении автоматов.}
%    \label{fig:product-automaton-example}
%\end{marginfigure}
%


%\section{Вопросы и задачи}
%
%Построить базу.
%
%Научиться выполнять запросы через линейку.

